\documentclass[12pt]{article}
\usepackage[utf8]{inputenc}
\usepackage[T1]{fontenc}
\usepackage{hyperref}
\title{Karen, Pwo Western Language Resources for K-12 Curriculum Development}
\author{Compiled by Claude Sonnet 5.5 (claude-sonnet-5-5) for Andr\'as Kornai}
\date{}
\begin{document}
\maketitle

\section{Introduction}
Western Pwo Karen (ISO 639-3 pwo; also Delta Pwo, Bassein Pwo, Phlong, Mon Karen) is spoken by about 210,000 people, mainly in the Ayeyarwady, Bago and Yangon regions of Myanmar, and is written in a Burmese-script-based orthography. Thai varieties of Pwo have separate codes (e.g. Kanchanaburi Pwo, kjp), so resources must be checked against the code. Little educational material was found; the dossier is short.

\section{Primary Educational Resources}
\begin{itemize}
\item Karen Language Resources (Minnesota Karen organisation resource page; mainly Karen languages generally, content for Pwo not verified): \url{https://mnkaren.org/resources/karen-language-resources/}
\end{itemize}
No Pwo-specific primers or textbooks were found in search results.

\section{Linguistic and Reference Materials}
\begin{itemize}
\item Omniglot, Western Pwo writing: \url{https://omniglot.com/writing/westernpwo.htm}
\item ASJP entry for Karen Pwo Western: \url{https://asjp.clld.org/languages/KAREN_PWO_WESTERN}
\item UNHCR news item on a committee beginning a English--Myanmar--Karen languages dictionary (completion not confirmed): \url{https://data.unhcr.org/en/news/20297}
\item Karen News tag page on Karen languages: \url{https://karennews.org/tag/karen-languages/}
\end{itemize}

\section{Religious and Community Resources}
\begin{itemize}
\item Asia Harvest profile of Pwo Karen Western (lists Bible 1883, New Testament 1860, portions 1845): \url{https://www.operation.asiaharvest.org/MYANMAR-Pwo-Karen-Western.pdf}
\end{itemize}

\section{Digital and Audio Resources}
\begin{itemize}
\item Global Recordings Network, Karen, Pwo Western: \url{https://globalrecordings.net/en/language/pwo}
\item Global Recordings Network, Karen, Pwo Western (alternate listing): \url{https://glorecnet-website.k.globalrecordings.net/en/language/416}
\item Global Recordings Network, Pwo Western in Thailand (Kanchanaburi Pwo, a different variety): \url{https://glorecnet-website.k.globalrecordings.net/en/language/11631}
\end{itemize}

\section{Curriculum Development Context}
The language is not known to be taught in schools. Burmese-script scripture and an unconfirmed planned dictionary are the main written bases found; any curriculum would need to build primers and readers largely from scratch, working with community organisations.

\end{document}